\section{Conclusion}

We began with an apparent tradeoff: converting transformers to sub-quadratic architectures discards few-shot adaptation capability. This paper shows the tradeoff is not fundamental---it stems from how conversion is performed, not from architectural constraints.

TC-SSM integrates task conditioning directly into the conversion process, preserving the adaptation manifold rather than attempting to recover it post-hoc. Our approach demonstrates that self-supervised clustering discovers meaningful task structure (purity 0.74), low-rank modulation adds negligible overhead (0.85$\times$ inference time), and integrated task conditioning preserves adaptation (0.25\% gap, 14 steps).

\subsection{Future Directions}

\textbf{Understanding the efficiency gain.} The sub-unity overhead was unexpected. Roofline analysis would provide definitive evidence for our memory bandwidth hypothesis.

\textbf{Addressing sample imbalance.} Balanced sampling during contrastive training could improve task embedding quality for rare task types.

\textbf{Scaling validation.} Experiments at 1B+ scale would confirm whether the mechanism transfers.

\textbf{Beyond classification.} Extending TC-SSM to generation tasks would establish broader applicability.

The efficiency-adaptation tradeoff has shaped how practitioners think about model deployment. Our work suggests this constraint can be engineered away through co-design of conversion and adaptation.
